\chapter{Intervolsorial Pediments and Marine Biomass}
\label{ch:20}

The marine counterpart of the terrestrial program of the preceding chapter couples the cultivation of kelp to the energy of the tide and to the processing of the harvest, and it is organized around structures that the author's earlier notes call volsorial and intervolsorial pediments. The descriptions given there are not mutually consistent in every detail, and the book uses the following reading of them. A volsorial pediment is a large tidally powered structure, described as a thousand meters tall, that serves as a gravitational battery: tidal energy is used to raise mass, the mass stores the energy as gravitational potential, and its descent returns the energy on demand. In oceanic versions the pediment is a platform that also carries kelp cultivation, with kelp-based parachutes and cables tethered by mycelial networks, and in some descriptions the platforms are mobile. Intervolsorial pediments are the suspended terraces or habitable levels between and among such structures, which the notes envision as housing on the order of two hundred thousand people each and as multifunctional sites for the recovery of resources, and which in the later orbital proposals serve as landing hubs for nested expandable towers. A further and more speculative idea, the hurricane mitt, would deploy an array of some ten thousand pediments with kelp parachutes to engage a tropical cyclone. The book treats the energy-storage function and the kelp coupling as the core of the concept, both of which can be analyzed with elementary mechanics, and treats the habitation and storm-engagement functions as conjectural elaborations whose assessment awaits specification. The definition used here is therefore a reading of scattered descriptions and is to be superseded by the author's fuller specification. The chapter analyzes three quantities: the energy of the tide that charges the battery, the storage that a gravitational battery of a stated height can hold, and the supply of nutrients that the tidal exchange delivers to the cultivation. The condition of marine and terrestrial biodiversity against which any such program would be judged is summarized in the global assessment of the intergovernmental platform~\cite{ipbes2019}.

The energy of a tidal lagoon is a matter of elementary mechanics. A basin of area $A$ that fills at high tide and is emptied through turbines at low tide, when the level outside has fallen by the tidal range $R$, releases the potential energy of the water that it holds above the outside level, and if the water is released as the level falls the available energy per tide is $\tfrac12\rho gAR^2$. The dependence on the square of the range explains the geographical concentration of tidal power: it is practicable where the range is several meters, which occurs in a small number of estuaries and bays of the world, and impracticable elsewhere. The best-known operating example, the barrage on the Rance estuary in France, encloses a basin of about $22\ \mathrm{km^2}$ in a region whose mean range is about eight meters, and a calculation with these figures gives about $7\times10^{12}$ J, or $1.97$ GWh, per tide, which if realized once in each tidal cycle of $12.42$ hours corresponds to a mean power of about $158$ MW. The actual mean output of the station, of the order of $500$ GWh per year, is about $57$ MW, roughly a third of that figure. The shortfall is a typical one, resulting from the partial use of the range, the efficiency of the turbines, the requirement of operating within a variable tide, and the need to hold water during parts of the cycle. A reasonable planning assumption is that an installation will deliver a quarter to a third of the energy that the formula gives for the full range, and the estimate of any proposal should be accompanied by the corresponding derating.

The coupling to cultivation depends on a different property of the exchange. Kelp, like all photosynthetic organisms, requires nutrients, among which nitrogen is commonly the limiting one in temperate coastal waters, and in an enclosed lagoon the supply of nitrogen is set by the rate at which the water is exchanged with the sea and by the concentration of the nutrient in the incoming water. Because the biomass takes up nutrient in proportion to its growth, the exchange rate places a ceiling on the production that the lagoon can sustain, whatever the density of the plants and the quantity of light. The formalization derives the ceiling: production per unit volume cannot exceed the product of the exchange rate and the external nutrient concentration, divided by the nutrient content of the biomass. For an exchange of $0.2$ volumes per day, an external concentration of $10$ millimoles of nitrogen per cubic meter, and a nitrogen content of two per cent of dry mass, the ceiling is $1.4$ grams of dry biomass per cubic meter per day, which for a lagoon five meters deep corresponds to about $26$ tonnes of dry mass per hectare per year. The figure is illustrative, but it indicates the magnitude: the nutrient supply of an ordinary coastal exchange can support production of the order that cultivated kelp achieves, and cultivation at greater density than the supply can bear would simply starve. The result also shows what the tidal energy and the cultivation have in common: both are proportional to the exchange, and both are maximized by a design that moves a large volume of water through the structure, so that the interests of the two functions are aligned in the dimension that matters most. Two kinds of precedent should be kept apart in reading what follows. The practical case for farming kelp together with shellfish in three-dimensional ocean plots, as a working livelihood and not as a laboratory proposal, is made from experience in Smith's account of regenerative ocean farming~\cite{smith2019}. The imaginative precedent, in which a living marine ecosystem is treated as part of the infrastructure of a civilization and not as a crop, belongs to the futurist and science-fiction tradition, of which Brin's novels are a leading example~\cite{brin1983}; such works supply a way of posing the question and no evidence about productivity, sequestration, or ecological consequence, which must come from the oceanographic literature.

The ecological assessment is governed by the standard of the author's earlier work, which is that the objective is diverse, geographically available, seasonally synchronized, and recoverable ecological production and not the maximum of biomass. Kelp cultivation can serve that objective when it is carried out as the cultivation of a known species in a defined place, with measurable effects on the food web, on the oxygen and carbon chemistry of the water, and on the organisms that depend on the habitat it creates. It can also fail the objective, through the escape of cultivated genotypes into wild populations, the transmission of disease, the shading of the benthos, and the concentration of grazers. The experience with the Sargassum belt in the Atlantic supplies the contrary case. Sargassum in the open ocean is productive habitat, but the same biomass stranded on a coast blocks light, consumes oxygen, and releases hydrogen sulfide and ammonia as it decomposes. The belt reached the highest abundance in the satellite record in 2025 and remained near that level afterward, with several factors, among them river nutrient discharge, atmospheric deposition, upwelling off West Africa, and changes in wind and current, implicated and none established as the sole cause. Biomass increased, and the function of the ecosystem in the places where it was deposited reversed. The lesson is that the quantity of biomass is not a measure of health, and that biomass the receiving system cannot incorporate is a burden and not a resource.

The processing of the harvest is the place where the argument returns to the energy balance. Kelp is wet: freshly harvested it is typically about nine-tenths water by mass. Any process that requires dry feedstock must first remove the water, and the latent heat of evaporation of water, $2.26$ MJ per kilogram, is large. For a tonne of kelp with eighty-eight per cent water the drying energy is about $1.99$ GJ, or $552$ kWh, whereas the dry matter, taken at a heating value of $12$ MJ per kilogram, contains $1.44$ GJ. The energy required to dry the feedstock thermally is therefore about $1.4$ times the energy it contains, and a conversion pathway that begins with thermal drying cannot be energetically self-sustaining unless the heat is obtained from a source other than the feedstock. This is the quantitative reason for the interest in processes that operate on wet material, among which hydrothermal conversion is the principal class. The Caldera Reactor, an oceanic thermopneumatic biomass-processing system developed in the author's earlier work, is a proposal of this kind. It is described as using geothermal steam to lift a compression plate, drawing seawater into the chamber by the vacuum produced as the steam condenses, compressing layered kelp, peat, and sediment, recovering energy through turbines, and yielding biocrude and feedstock for biochemicals and bioplastics, with the conversion handled by hydrous pyrolysis for wet feedstock and anhydrous pyrolysis for drier streams. It is a system whose individual elements are plausible and whose integration has not been demonstrated at scale, and it is correctly assigned to the second readiness class. The author's design rule that cultivated kelp should be the normal feedstock, with excess Sargassum intercepted before it reaches the shore as a secondary and disturbance-responsive input, is a sound one. A processing system that depended structurally on the bloom would have a financial interest in the continuation of a pathology, which is the situation of any enterprise that profits from the damage it is presented as remedying.

Finally, the program is subject to the constraints of governance that apply to any use of the coast. Coastal waters are used for fishing, navigation, recreation, conservation, and the cultural practices of coastal communities, many of whom have claims of long standing. A structure that encloses a lagoon excludes some of these uses and alters others, and the effects on migrating fish and on sediment transport along the coast can extend far beyond the structure. The assessment must therefore include the effect on sediment budgets and on the fisheries of the region, the provision of passage for migratory species, and the mechanism by which the affected communities participate in decisions. The experience of tidal barrages, which have altered the ecology of the estuaries in which they were built and are nonetheless regarded by some as successes, indicates that the effects can be substantial and the judgments contested.

\section*{Formalization}

Three results are given, on the energy of a lagoon, the nutrient ceiling on production, and the energetics of drying.

For a lagoon of horizontal area $A$ (taken constant over the range), water density $\rho$, and a tidal range $R$, suppose the lagoon fills at high tide, and then is emptied through turbines while the outside level falls from the high-tide level to the low-tide level, with generation continuing until the lagoon level reaches the outside level. In the idealization in which the lagoon is emptied at the moment of low tide and all potential energy of the retained water is converted, the energy per tide is the work to lift the retained volume $AR$ by its mean height $R/2$.

\begin{proposition}
The energy released per tide by complete emptying is $E_{\mathrm{tide}}=\tfrac12\rho gAR^2$. If the lagoon is emptied once in each tidal period $T$, the mean power in the idealized cycle is $\bar P=\tfrac12\rho gAR^2/T$. In practice the delivered energy is $\eta E_{\mathrm{tide}}$ with $\eta<1$ accounting for the partial use of the range, the turbine efficiency, and constraints of operation.
\end{proposition}

\begin{proof}
The mass of water retained above the low-tide level is $\rho AR$, whose centre of mass is at height $R/2$ above that level. The potential energy is $\rho AR\cdot g\cdot R/2$. With one emptying per semidiurnal tidal period $T=12.42$ hours, the mean power is $E_{\mathrm{tide}}/T$.
\end{proof}

With $\rho=1025\ \mathrm{kg\,m^{-3}}$, $A=22\ \mathrm{km^2}$, and $R=8$ m, $E_{\mathrm{tide}}=7.08\times10^{12}$ J $=1.97$ GWh, and with one emptying per period $T=12.42$ h $=4.47\times10^4$ s the mean power is $158$ MW. A mean output of $57$ MW corresponds to $\eta\approx0.36$.

For the gravitational battery, a mass $M$ raised through a height $H$ stores $E=MgH$, with round-trip efficiency $\eta_{rt}$. The energy per tonne at $H=1000$ m is $9.81\times10^{6}$ J, which is $2.725$ kWh, so that the mass required to store an energy $E_{\mathrm{st}}$ is $M=E_{\mathrm{st}}/(gH\eta_{rt})$. A settlement of $2\times10^5$ persons with a mean demand of $1$ kW per person, which is $200$ MW, requires $4.8$ GWh to bridge a full day, which at $H=1000$ m and $\eta_{rt}=1$ is $1.76\times10^6$ tonnes, comparable in order of magnitude to the mass of the largest masonry monument, and at $\eta_{rt}=0.8$ is $2.2\times10^6$ tonnes. For comparison, at a representative $200$ kWh per tonne for a lithium-ion pack the same energy would be stored in $2.4\times10^4$ tonnes, about $1\%$ of the gravitational mass, so that the gravitational battery trades a vastly larger mass for the absence of chemical degradation and for the use of inexpensive material, a trade whose merit is an engineering and economic question.

For the nutrient ceiling, let the lagoon have volume $V$, exchange rate $e$ (volumes per unit time), external nutrient concentration $N_{\mathrm{ext}}$, internal concentration $N$, and uptake of $y$ units of nutrient per unit of biomass produced. Let $P$ denote production of biomass per unit volume per unit time.

\begin{proposition}
At steady state, $e(N_{\mathrm{ext}}-N)=yP$. Hence $P\le eN_{\mathrm{ext}}/y$, with equality if and only if the internal nutrient concentration is driven to zero. The ceiling is proportional to the exchange rate $e$ and cannot be raised by increasing the density or the light supply.
\end{proposition}

\begin{proof}
A nutrient balance on a well-mixed lagoon gives $V\dot N=eV(N_{\mathrm{ext}}-N)-yPV$, and at steady state the bracket vanishes. Since $N\ge0$, $yP=e(N_{\mathrm{ext}}-N)\le eN_{\mathrm{ext}}$.
\end{proof}

With $e=0.2\ \mathrm{d^{-1}}$, $N_{\mathrm{ext}}=10\ \mathrm{mmol\,m^{-3}}$, and $y=0.02/14\ \mathrm{mol\,g^{-1}}=1.43\ \mathrm{mmol\,g^{-1}}$ of dry mass, the ceiling is $P=1.4\ \mathrm{g\,m^{-3}\,d^{-1}}$. Over a depth of $5$ m this is $7\ \mathrm{g\,m^{-2}\,d^{-1}}$, or $25.6\ \mathrm{t\,ha^{-1}\,yr^{-1}}$ of dry mass. The corresponding sustainable harvest policy is obtained from a logistic growth law with harvest rate $h$ and effective growth rate $r_{\mathrm{eff}}$ (the rate of growth under the nutrient conditions): the equilibrium biomass is $K^\ast=K_{\max}(1-h/r_{\mathrm{eff}})$ and the yield $hK^\ast$ is maximal at $h=r_{\mathrm{eff}}/2$, where it equals $r_{\mathrm{eff}}K_{\max}/4$. The maximum-yield policy sits at half the carrying capacity and is the boundary of the region in which a perturbation of $r_{\mathrm{eff}}$ can cause collapse, so that a policy chosen with margin sets $h$ below $r_{\mathrm{eff}}/2$.

For the drying energetics, let wet biomass of mass $m$ have a water mass fraction $w$ and let the dry matter have heating value $q$ per unit mass. The thermal drying energy is $\ell_v wm$, where $\ell_v=2.26$ MJ kg$^{-1}$, and the energy content is $q(1-w)m$.

\begin{proposition}
The ratio of drying energy to energy content is $\mathcal R=\ell_vw/(q(1-w))$. It exceeds $1$ when $w>q/(q+\ell_v)$, a threshold of $0.842$ for $q=12$ MJ kg$^{-1}$ and $0.869$ for $q=15$ MJ kg$^{-1}$, and it increases without bound as $w\to1$. For kelp at $w=0.88$ and $q=12$, $\mathcal R=1.38$.
\end{proposition}

\begin{proof}
Direct computation: $\mathcal R>1$ iff $\ell_vw>q(1-w)$, i.e. $w(\ell_v+q)>q$. For $w=0.88$, $\mathcal R=2.26\cdot0.88/(12\cdot0.12)=1.381$.
\end{proof}

The result shows why wet-processing pathways are a requirement and not an optimization for feedstocks of high moisture content, and why the heating value is of secondary importance: even at $q=15$ MJ kg$^{-1}$ the ratio is $1.10$.
